% Source-aligned English for printed p.126.
\SourcePage{126}

\LayerTarget{tohoku:section:1.3}
\noindent\textbf{1.3. Additive categories.}
An \textit{additive category} is a category $C$ for which, for every pair
$(A,B)$ of objects of $C$, an abelian-group law is given on
$\mathrm{Hom}(A,B)$, in such a way that composition of morphisms is a bilinear
operation.

\LayerTarget{tohoku:section:1.3:p1}
We assume moreover that the sum and the product of any two objects $A,B$ of
$C$ exist. It is in fact enough to postulate the existence of the sum
\textit{or} the product of $A$ and $B$; the existence of the other then follows
easily, and moreover $A+B$ is canonically isomorphic to $A\times B$. (Assuming,
for example, that $A\times B$ exists, consider the morphisms
$A\longrightarrow A\times B$ and $B\longrightarrow A\times B$ whose
components are $(I_A,0)$ and $(0,I_B)$, respectively; one verifies that this
gives a representation of $A\times B$ as a direct sum of $A$ and $B$.)
Finally, we postulate the existence of an object $A$ such that $I_A=0$; it is
called a \textit{null object}, or a \textit{zero}, of $C$. Equivalently,
$\mathrm{Hom}(A,A)$ is reduced to zero, or, again, for every $B\in C$,
$\mathrm{Hom}(A,B)$ (or $\mathrm{Hom}(B,A)$) is reduced to zero. If $A$ and
$A'$ are two null objects, there is then a unique isomorphism from $A$ to $A'$
(namely the unique element $0$ of $\mathrm{Hom}(A,A')$!); we shall therefore
identify all the null objects of $C$ with a single one, denoted $0$, by an
abuse of notation.

\LayerTarget{tohoku:section:1.3:p2}
The dual category of an additive category is again an additive category.

\LayerTarget{tohoku:section:1.3:p3}
Let $C$ be an additive category and let $u:A\longrightarrow B$ be a morphism
in $C$. For $u$ to be injective (respectively, surjective), it is necessary
and sufficient that there be no nonzero morphism whose composition on the
right (respectively, on the left) with $u$ is zero. A \textit{generalized
kernel} of $u$ is a monomorphism $i:A'\longrightarrow A$ such that the
morphisms $C\longrightarrow A$ that are right zero-divisors for $u$ are
exactly those that factor as

\[
  C\longrightarrow A'\mathrel{\mathop{\longrightarrow}^{i}}A.
\]

\noindent Such a monomorphism is defined up to equivalence (cf.
No.~1); hence, among the generalized kernels of $u$ (if any exist), there is
exactly one that is a subobject of $A$: it is called \textit{the kernel of
$u$}, and is denoted $\mathrm{Ker}\,u$. Dually, one defines the
\textit{cokernel of $u$} (which is a quotient object of $B$, if it exists),
denoted $\mathrm{Coker}\,u$. The \textit{image} (respectively,
\textit{coimage}) of the morphism $u$ is the kernel of its cokernel
(respectively, the cokernel of its kernel), if it exists; it is thus a
subobject of $B$ (respectively, a quotient object of $A$)%
\LayerTarget{tohoku:section:1.3:footnote:1ter}\textsuperscript{1 ter)} and is
denoted $\mathrm{Im}\,u$ (respectively, $\mathrm{Coim}\,u$). If $u$ admits an
image and a coimage, there is a unique morphism
$\bar u:\mathrm{Coim}\,u\longrightarrow\mathrm{Im}\,u$ such that $u$ is
identical to the composite

\[
 A\longrightarrow\mathrm{Coim}\,u\longrightarrow
 \mathrm{Im}\,u\longrightarrow B,
\]

\noindent where the outer morphisms are the canonical morphisms.

\LayerTarget{tohoku:section:1.3:p4}
A functor $F$ from an additive category $C$ to another such category $C'$ is
called an \textit{additive functor} if, for two morphisms
$u,v:A\longrightarrow B$ in $C$, one has
$F(u+v)=F(u)+F(v)$. The analogous definition applies to multifunctors.
Composites of additive functors are additive. If $F$ is an additive functor,
$F$ carries a finite direct sum of objects $A_i$ to the direct sum of the
$F(A_i)$.

\par\medskip\hrule\smallskip
\LayerTarget{tohoku:section:1.3:footnote:1ter:text}%
\noindent
\textsuperscript{1 ter)} A more natural definition of the image of $u$, to
tell the truth, would be to take the smallest subobject $B'$ of $B$ (if one
exists) such that $u$ comes from a morphism from $A$ to $B'$. This definition
is equivalent to the one given in the text only when $C$ is an abelian
category (cf. 1.4.)

\TohokuContinuation{127}
